\chapter{Undo and Time}
\label{ch:undo}

Every transformation covered in Chapters~\ref{ch:operators} and~\ref{ch:insert} is reversible, and this chapter is concerned with that reversal: not as an emergency escape hatch, though it functions as one, but as a navigable structure in its own right, one with more depth to it than the simple linear "step backward" model a reader arriving from other editors might expect.

\section{Basic Undo and Redo}

\texttt{u} undoes the most recent change, and can be pressed repeatedly to step back through successive earlier changes, one at a time. \texttt{Ctrl-r} redoes a change that was just undone, stepping forward again. \texttt{U} is a distinct, older command that undoes all changes made to the current line since it was last visited, which behaves unlike a simple repeated \texttt{u} in one specific respect worth flagging: pressing \texttt{U} a second time undoes the effect of the first \texttt{U}, restoring the line to the state \texttt{U} had just undone it from, rather than continuing to undo further back, which makes \texttt{U} function as a toggle on a single line's state rather than as a general step-backward command.

What counts as a single "change" for the purposes of \texttt{u} is worth stating explicitly, since it is not always the same granularity a reader might assume. An entire Insert-mode session, from the moment Insert mode is entered until \texttt{escape} returns to Normal mode, counts as one change, however much text was typed during it; a single \texttt{u} after a long paragraph has been typed undoes the entire paragraph at once, not the most recently typed character. A single operator-motion command, such as \texttt{daw}, also counts as one change. This granularity is generally the intuitive one in practice, since it corresponds to what a user would describe as "one edit" in ordinary language, but it is worth knowing explicitly for the cases where it matters, such as a very long Insert-mode session that the user may want to partially, rather than entirely, undo.

\section{The Undo Tree}

The model implied by "step backward, step forward" is linear, and it is the model most editors implement: a single sequence of states, with undo and redo moving a pointer backward and forward along it. Vim's undo history is not linear; it is a tree. If the user undoes several changes and then makes a new edit rather than redoing, the changes that were undone are not discarded, as they would be in a strictly linear model. They remain in the undo history as a separate branch, reachable again later, while the new edit begins a second branch from the same point.

This matters concretely in a scenario every long editing session eventually produces: a paragraph is deleted, several unrelated edits are made elsewhere in the file, and only then does the user realize the original deletion should not have happened. In a linear undo model, undoing back to before the deletion would also undo all the unrelated edits made afterward, forcing a choice between recovering the paragraph and keeping the later work. Vim's tree-structured undo history does not force this choice, because the later edits and the deleted paragraph live on different branches of the same tree, both still present, both still reachable.

\texttt{g-} and \texttt{g+} move backward and forward through the undo history by \textit{time} rather than by tree position, meaning they follow the sequence in which changes actually occurred, across branches, rather than following the parent-child structure \texttt{u} and \texttt{Ctrl-r} follow within a single branch. \texttt{:undolist} displays the tree's branch points directly, each labeled with a number, and \texttt{:undo n} jumps directly to the state after change number \textit{n}, which is the most direct way to navigate to a specific branch once \texttt{:undolist} has been consulted to identify it.

\section{Persistent Undo}

By default, a buffer's undo history exists only for the duration of the editing session; closing the file and reopening it later begins with no undo history at all, even if the file was previously edited and saved. \texttt{set undofile}, placed in the vimrc introduced in Chapter~\ref{ch:installing}, changes this: Vim writes the undo history to a separate file on disk (by default alongside the edited file, or in a centralized directory if \texttt{undodir} is set), and reloads that history automatically the next time the same file is opened, meaning \texttt{u} can undo changes made in a previous, entirely separate editing session, potentially days or weeks earlier, exactly as though that session had never ended. This is one of the more consequential single-line additions to a minimal vimrc, and worth adding immediately once it is known to exist, since its absence is only ever noticed at the moment it would have been useful and no longer can be.

\section{Swap Files and Crash Recovery}

Separately from the undo history, Vim maintains a swap file for each buffer currently being edited (typically named \texttt{.filename.swp} in the same directory as the file, unless \texttt{directory} is configured otherwise), recording changes as they are made so that an abnormal termination, a crash, a lost SSH connection, a killed process, does not lose the in-progress edits entirely. On a normal, clean exit, the swap file is deleted automatically; its continued presence when a file is reopened is itself the signal that a previous session did not end cleanly, and Vim will warn the user of this directly rather than silently ignoring it.

When that warning appears, \texttt{vim -r} \textit{filename} recovers the buffer's state from the swap file as it stood at the last recorded point before the abnormal termination, rather than opening the file fresh from its last saved state on disk, which may be considerably further behind. Recovery should generally be followed by writing the recovered buffer to disk under a new name first, inspecting it, and only then deciding whether to overwrite the original file, since a swap file's contents, while usually reliable, are not guaranteed to reflect every last keystroke before the crash occurred.

\section{Undo as Confidence}

The material in this chapter is not, by itself, part of the operator-motion grammar Part II established, and it does not multiply against operators or text objects the way the material in Chapters~\ref{ch:operators} and~\ref{ch:insert} does. It is included in Part III regardless, immediately after the transformations it exists to reverse, because its practical effect on how the rest of this book should be read is significant: every operator, every insertion, every substitution covered up to this point is safe to attempt experimentally, on real files, because \texttt{u} is always available, the undo tree never discards a state the user has not explicitly abandoned, and persistent undo, once enabled, extends that safety net across sessions rather than confining it to the current one. A reader who has internalized this chapter should feel, for the remainder of the book, considerably less hesitant about trying an unfamiliar command on a real file than they otherwise might.
